% ===================================================================== Chapter 17
\chapter{Engineering Discussion}\label{ch:discussion}

\section{Flow Behavior}
The flow is a developing turbulent pipe flow with strong heating. Two features distinguish it from a textbook isothermal case. First, the
air density falls by about a fifth along the duct, so the flow accelerates and a third of the static pressure drop (149~Pa of 438~Pa) goes into
acceleration rather than friction. Second, heating lowers the wall friction relative to the constant-property correlation by about 9\,\%. The
second effect explains why the friction factor misses the Petukhov value by slightly more than 10\,\%, and it is a property of the physics, not
of the mesh. For the structural question, the flow matters mainly through the heat-transfer coefficient it produces.

\section{Thermal Behavior}
The thermal behaviour is controlled by the air-side film. The metal carries only a few percent of the thermal resistance, so the wall is
almost isothermal through its thickness (7.6~K across 10~mm at mid-span) but runs about 186~K above the bulk air near the outlet. Along the duct
the wall temperature follows the bulk temperature plus the film temperature difference, modified at both ends by conduction: the inlet end is
cooled by the high heat-transfer coefficient of the thermal entrance, which draws heat axially through the wall, and the outlet end is
warmed by the adiabatic end face. This is why the CFD maximum is 19~K below the one-dimensional estimate and why the mean temperature rise, which
drives the structural loads, is about 11\,\% lower than in the analytical model. The conjugate CFD was therefore necessary: the one-dimensional
model overestimates the structural loads because it cannot represent the entrance and end effects.

\section{Structural Behavior}
The structural behaviour is governed by restraint rather than by the gradient. Free expansion produces a growth of 1.84~mm but only 24~MPa of
stress, confined to the inlet end; restrained expansion stores the same strain as a nearly uniform compression of 582~MPa. The thermal
problem therefore translates into a restraint problem: the stress level depends on how much of the growth the mounting prevents, not on
the temperature distribution within the wall. The static state is elastic in every case, with the peak at the inlet-face edge where the
compression and the inlet gradient combine.

\section{Buckling Behavior}
The restrained duct is a stocky column carrying a thermal, displacement-controlled load close to its idealised Euler load. Under S1 the first
bifurcation appears at 1.108 times the applied load, before first yield. Two points qualify this result. The load is thermal, so a small
change in mean temperature (about 24~K) changes the load factor by the full 10.8\,\%; the stability statement is therefore as uncertain as the
mean wall temperature, which itself differs by 30~K between the analytical and CFD models. And the linear eigenvalue is an upper-bound-type
indicator for a perfect column; imperfections and inelastic behaviour act only to lower the real capacity. The additional LS-DYNA
analysis (Section~\ref{sec:lsdyna}) is consistent with this: with mode-shaped imperfections the lateral deflection and stress grow well below
$\lambda_1$, while the Southwell characteristic-load estimates stay within 1\,\% of the linear critical load.

\section{Effect of Velocity}
Velocity primarily affects cooling, Reynolds number and pressure drop. It changes the heat-transfer coefficient and the air temperature rise,
and through them the wall temperature; the structural response then follows the temperature. A 10\,\% lower velocity raises the wall
temperatures by about 25~K, the restrained force by 9.7\,\% and lowers $\lambda_1$ to 1.003, while the critical load itself hardly changes. The
pressure-drop penalty of higher velocity (+14.1\,\% for +10\,\% velocity) is the price of the cooler wall.

\section{Effect of Heat Flux}
Heat flux directly drives temperature and thermal stress. It acts on the wall temperature rise almost in proportion (slightly more than
proportionally), and every thermal-structural response follows: $\pm$10\,\% flux gives about $\pm$11\,\% restrained force and LC2 stress and moves
$\lambda_1$ from 1.255 to 0.988. It is the variable to which $\lambda_1$ is most sensitive. The pressure drop changes only through the gas
temperature.

\section{Effect of Wall Thickness}
Wall thickness strongly affects stiffness and buckling capacity while having relatively little thermal effect under the constant-total-heat-input
definition. Temperatures change by less than 1~K and the stress level by less than 0.5\,\%, but the end force scales with the area and the
critical load with the bending stiffness. A thicker wall carries a larger force but gains more capacity, because the radius of gyration
grows; a thinner wall loses capacity faster than it loses force. Thickness is therefore a structural design variable, not a thermal one, in
this range. With the heat flux held constant instead of the total heat input, the thermal effect would differ; that alternative was not studied.

\section{Effect of Support Condition}
The support condition dominates the idealised buckling factor. With identical static stress, the three idealised supports change $\lambda_1$ by
a factor of 3.9, far more than any operating or design variable in the study, and they decide whether bifurcation or first yield is reached
first. The mesh resolution, by comparison, changes $\lambda_1$ by less than $10^{-4}$ on the structural side and by up to about 2\,\% through the
CFD temperature field. The mesh is not the dominant uncertainty; the boundary condition is.

\section{Physical Interpretation}
Taken together, the results describe a chain in which each link amplifies the importance of the next. The air film sets the wall
temperature; the restraint converts that temperature into a large compressive force; and the end fixity decides whether that force is
stable. Pressure loading is negligible throughout ($7.4\times10^{-6}$\,\% of the peak stress). The engineering implication, within the
idealised model, is that the design question for such a duct is less about reducing the through-wall gradient than about how
much axial growth the mounting prevents and how the ends are restrained laterally and rotationally. Allowing controlled axial growth, or
providing defined lateral end restraint, are the two routes the model identifies; neither was analysed as a design and both would need the
real mounting to be defined.
